\documentclass[letterpaper]{article}

\usepackage[letterpaper,margin=1in]{geometry}

% Compile from the repository root so that this path resolves to
% the pinned WSP submodule.
\input{wsp/testing/test-case-library.tex}

\begin{document}

\def\TCID{%
  TC-0007
}

\def\TCTitle{%
  ListView rules
}

\def\TCPurpose{%
  Verify the rules of a list: the valid views, the layout of every view, hit
  testing, sorting, FindItem matching, SubItems limits, selection and label
  editing.
}

\def\TCPriority{%
  \TCPriorityHigh
}

\def\TCRequirementRef{%
  REQ-0014, REQ-0031, REQ-0033, REQ-0034
}

\def\TCDesignRef{%
  Not applicable; the rules are in the named header.
}

\def\TCFeatureRef{%
  \texttt{include/wfc/mscomctl/list.hpp} and \texttt{tests/list\_test.cpp} and
  the runner \texttt{tests/tc-0007-list.ps1}
}

\def\TCTestTechnique{%
  Equivalence partitioning and boundary value analysis of the rule functions
  (ISO/IEC/IEEE 29119-4 specification-based techniques).
}

\def\TCPreconditions{%
  The library is built standalone (\texttt{cmake -S . -B out/build}, then
  \texttt{cmake --build out/build --config Debug}) with
  \texttt{WFC\_BUILD\_TESTS} on; CTest runs the test as \texttt{TC-0007-list}.
}

\def\TCEnvironment{%
  Any target the library builds for (Windows x86, x64 and ARM64 with MSVC,
  Linux through Winelib); local runs use the x64 Debug configuration of the
  standalone build. The test needs no window, registry or network.
}

\def\TCAssumptions{%
  The rule header under test is the one built into the library; the test is
  compiled from the same source revision.
}

\def\TCInputData{%
  Views -1 to 4, lists of 1 to 5 items in each view and size, sort keys -3 to
  3, the texts \texttt{alpha}, \texttt{Bravo}, \texttt{Same}, and selection
  vectors.
}

\def\TCInitialState{%
  A fresh process; the test keeps no state between runs and holds no files or
  handles.
}

\def\TCProcedure{%
  \begin{enumerate}
    \item Run \texttt{mscomctl\_list\_test} (CTest \texttt{TC-0007-list}).
    \item Check each view number, then lay out items in Report (with and
      without headers), Icon and List views and hit-test them.
    \item Check sort sources and ascending and descending comparison, whole
      and partial FindItem matches, and the sub-item read and write checks.
    \item Select items with and without Ctrl, and ask the second-click, extend
      and label-edit questions.
  \end{enumerate}
}

\def\TCExpectedResult{%
  \begin{enumerate}
    \item Views 0 to 3 are valid.
    \item Report rows sit below the header (240 twips) and at the top without
      one; icons wrap at the width; the list view fills columns downwards; an
      unset height still lays out.
    \item Hit testing finds the row and nothing right of or below the rows.
    \item Sort by the item or by a sub item, ascending or descending, ignoring
      case, equal texts in neither order; a whole match ignores case, a
      partial match is a prefix and the empty text matches.
    \item Reading sub items beyond the count is 35600, writing a negative
      index is 35600, writing grows the sub items up to the index.
    \item A click replaces the selection, Ctrl toggles and extends only with
      MultiSelect; automatic editing starts on a second click; a scripted edit
      starts unless the style is manual; bytes compare as unsigned.
    \item The program prints \texttt{list: ok} and exits with code 0.
  \end{enumerate}
}

\def\TCPostConditions{%
  The process ends; nothing is written outside CTest's own log.
}

\TCTemplate

\end{document}
